\section{Related Work}
\label{sec:related}

\paragraph{Unsupervised environment design.}
PAIRED~\citep{dennis2020paired} frames curriculum design as a game between
a learned level generator and two students, a protagonist and an
antagonist, whose return gap approximates regret.
PLR~\citep{jiang2021plr} keeps a buffer of sampled levels prioritised by a
score such as positive value loss and replays them by rank and staleness;
\plr~\citep{jiang2021replay} trains only on replayed levels, which makes it
a regret-minimax method.
ACCEL~\citep{parkerholder2022accel} extends \plr\ by editing replayed
levels and scoring the edits as new candidates, with a generator of simple
levels, and reports gains on 2D BipedalWalker terrain.
SFL~\citep{rutherford2024noregrets} shows that common regret approximations
track success rate rather than regret, scores levels by learnability
$p(1-p)$ for success probability $p$, and proposes an adversarial
CVaR-style evaluation that we adapt.
minimax~\citep{jiang2023minimax} and JaxUED~\citep{coward2024jaxued}
provide accelerated reference implementations on 2D domains.
Our benchmark keeps these algorithms close to their published form and
changes the domain: levels are heightfields for a 3D quadruped.

\paragraph{Terrain curricula for legged locomotion.}
\citet{lee2020learning} train ANYmal over parameterised terrain with an
adaptive curriculum that samples terrain parameters of intermediate
difficulty. \citet{rudin2022learning} introduce a game-inspired curriculum
in which robots move to harder terrain rows after walking far enough; it is
the default in massively parallel locomotion training.
LP-ACRL~\citep{li2026lpacrl} samples tasks by estimated learning progress
and reaches high-speed ANYmal~D locomotion across stairs, slopes, gravel and
low-friction ground.
TACL~\citep{liu2026tacl} learns a difficulty function over
trajectory-level tasks on unstructured terrain maps and proposes the next
tasks from it.
PSC~\citep{ovinnikov2026psc} prioritises terrain contexts and randomised
events by a distributional safety critic and compares against terrain
progression, advantage-based replay and learning-progress curricula.
Eurekaverse~\citep{liang2024eurekaverse} uses an LLM to write progressively
harder parkour courses for a quadruped.
\citet{zhang2022terrain} generate terrain datasets by authoring and active
learning as a prototype robustness benchmark.
Each of these works proposes one curriculum or dataset and compares it
with a few baselines; none runs the standard UED family (PAIRED, \plr,
ACCEL, SFL) side by side under a shared level space and evaluation.
\todo{Final literature sweep before submission (2025--2026 UED and
locomotion-curriculum papers); check whether TACL or PSC include PLR/ACCEL
baselines in their experiments, not only in related work.}
